\documentclass[border=6pt]{standalone}
\usepackage[T1]{fontenc}
\usepackage{lmodern}
\usepackage{amsmath}
\usepackage{tikz}
\usetikzlibrary{arrows.meta,decorations.pathreplacing}
\tikzset{eleg/.style={dash pattern=on 2.6pt off 2pt}}
\usepackage{pgfplots}
\pgfplotsset{compat=1.18}
\usepgfplotslibrary{colormaps}
\pgfplotsset{colormap={mplhot}{rgb255=(11,0,0) rgb255=(32,0,0) rgb255=(53,0,0) rgb255=(74,0,0) rgb255=(97,0,0) rgb255=(118,0,0) rgb255=(139,0,0) rgb255=(160,0,0) rgb255=(184,0,0) rgb255=(205,0,0) rgb255=(226,0,0) rgb255=(247,0,0) rgb255=(255,16,0) rgb255=(255,37,0) rgb255=(255,58,0) rgb255=(255,79,0) rgb255=(255,102,0) rgb255=(255,123,0) rgb255=(255,144,0) rgb255=(255,165,0) rgb255=(255,189,0) rgb255=(255,210,0) rgb255=(255,231,0) rgb255=(255,252,0) rgb255=(255,255,31) rgb255=(255,255,62) rgb255=(255,255,94) rgb255=(255,255,125) rgb255=(255,255,160) rgb255=(255,255,192) rgb255=(255,255,223) rgb255=(255,255,255)}}
\definecolor{inkPrim}{HTML}{1B1B1B}
\definecolor{coral}{HTML}{F0653F}

\begin{document}
\begin{tikzpicture}[font=\rmfamily]
% ===================== dynamical structure factor S(q,omega), L=12 =====================
\begin{axis}[
  name=main,
  width=13.8cm, height=8.6cm, enlargelimits=false, axis on top,
  colormap name=mplhot, point meta min=0, point meta max=1.33,
  xmin=0, xmax=2, ymin=0, ymax=12,
  xtick={0,0.5,1,1.5,2}, xticklabels={$0$,$\tfrac12$,$1$,$\tfrac32$,$2$},
  ytick={0,2,4,6,8,10,12},
  minor x tick num=1, minor y tick num=1,
  xmajorgrids, ymajorgrids, xminorgrids, yminorgrids,
  major grid style={white, opacity=0.20, line width=0.3pt},
  minor grid style={white, opacity=0.09, line width=0.2pt},
  % 4.968759 = Delta from data/charge_gap.json (one-particle channel, NOT measured in S(q,w));
  % 8.779 = peak(q/pi=1.0) from sqw_edges.dat. Both verified against deposited files.
  extra y ticks={4.968759,8.779},
  extra y tick labels={\textcolor{coral}{$\mathbf{\approx5}$},\textcolor{coral}{$\mathbf{8.8}$}},
  extra y tick style={grid=none, major tick length=7pt, tick style={coral, line width=1.3pt},
    tick label style={font=\footnotesize}},
  extra x ticks={1.0}, extra x tick labels={\textcolor{coral}{$\mathbf{2k_F}$}},
  extra x tick style={grid=none, major tick length=6pt, tick style={coral, line width=1.2pt},
    tick label style={font=\footnotesize, yshift=-9pt}},
  tick align=outside, ytick pos=left,
  axis line style={inkPrim, line width=0.6pt}, tick style={inkPrim, line width=0.6pt},
  tick label style={font=\footnotesize},
  xlabel={momentum transfer \; $q/\pi$}, ylabel={frequency \; $\omega/t$}, label style={font=\small},
  colorbar, colorbar style={width=0.28cm, ylabel={response intensity \; $S(q,\omega)$},
    ylabel style={font=\small}, yticklabel style={font=\footnotesize}, ytick={0,0.4,0.8,1.2}},
]
  \addplot graphics[xmin=0,xmax=2,ymin=0,ymax=12]{sqw_field.png};

  % ---- parameters (identical style to A(k,omega)) ----
  \node[white, font=\scriptsize, fill=black, fill opacity=0.5, text opacity=1, inner sep=1.6pt,
        rounded corners=1pt, anchor=north east] at (axis cs:1.98,11.7) {$U/t=8,\ \eta=0.18\,t$};

  % ---- 2k_F guide line (same white dash-dot as k_F in A(k,omega)) ----
  % 8.779 = peak(q/pi=1.0000) from release/paper/figs/sqw_edges.dat, regenerated by
  % src/make_sqw_edges.py from data/sqw_L12.json. Hand-placed but VERIFIED against the file.
  \draw[white, dash pattern=on 5pt off 2.5pt on 1.5pt off 2.5pt, line width=0.7pt]
        (axis cs:1.0,0) -- (axis cs:1.0,8.779);

  % ---- the collective feature (title, same style as the Hubbard-band labels) ----
  \node[white, font=\footnotesize\itshape, fill=black, fill opacity=0.55, text opacity=1, inner sep=2pt,
        rounded corners=1.5pt, anchor=north west] at (axis cs:0.02,11.8) {charge (holon--antiholon) band};

  % ---- lower edge (charge threshold): white densely dotted, same style as the spinon curve ----
  %
  % PROVENANCE, 2026-09-18. This plot previously read
  %     rebuild/figs/sqw_edges.dat
  % which lives outside the deposit and still contains two FABRICATED rows,
  %     0.0000 4.969 5.000 5.000   and its mirror at qpi=2.0000,
  % produced by no script in the repository. q=0 is inaccessible for this observable by
  % construction -- src/sqw_lanczos.py:30 skips it and src/sqw_field.py:3 states
  % S(q=0,w>0)=0 by density conservation -- and 5.000 is not a point of the frequency grid
  % (neighbours 4.9691 and 4.9925). The value 4.969 is Delta = mu+ - mu- taken from the
  % ONE-PARTICLE channel, not a measurement of S(q,w). Drawing the edge through it made the
  % caption's claim ("as q->0 its lower edge approaches Delta") circular: the point that
  % demonstrated the convergence had been set by hand equal to the thing it demonstrated.
  % The plot now reads the regenerated, deposited file, whose 11 rows are reproduced byte for
  % byte by src/make_sqw_edges.py from data/sqw_L12.json. The lowest accessible datum is
  % onset(q/pi = 1/6) = 5.600.
  \addplot[white, line width=0.7pt, densely dotted, forget plot, smooth]
        table[x=qpi,y=onset]{sqw_edges.dat};
  \draw[white, line width=0.5pt, dotted] (axis cs:1.40,3.9) -- (axis cs:1.50,5.55);
  \node[white, font=\scriptsize, fill=black, fill opacity=0.6, text opacity=1, inner sep=1.5pt,
        rounded corners=1pt, anchor=west] at (axis cs:1.06,3.5) {onset: lowest pole at each $q$};

  % ---- effect 1: the continuum lies above the Mott gap Delta = mu+ - mu- ----
  % Delta is a real, independently computed quantity -- src/charge_gap_ed.py -> data/charge_gap.json,
  % exact diagonalization at L=12, U/t=8: mu+ = 6.484380, mu- = 1.515620, Delta = 4.968759 t --
  % so marking it is legitimate. What is NOT legitimate is marking it as a POINT OF THIS CURVE at
  % an inaccessible momentum. It is therefore drawn as what it is: a horizontal reference level
  % imported from the one-particle channel, spanning the panel, with no data point on it.
  \draw[inkPrim, line width=0.7pt, dash pattern=on 3pt off 2pt]
        (axis cs:0.00,4.968759) -- (axis cs:2.00,4.968759);
  \node[white, font=\footnotesize\bfseries, fill=black, fill opacity=0.80, text opacity=1,
        inner sep=2.5pt, rounded corners=2pt, anchor=west] at (axis cs:0.30,2.35)
        {charge gap $\Delta=\mu^{+}\!-\mu^{-}=4.97\,t$ \; (from $A(k,\omega)$, not measured here)};
  \draw[white, line width=0.7pt, -{Stealth[length=1.8mm]}] (axis cs:0.42,2.85) -- (axis cs:0.30,4.80);

  % ---- effect 2: the mode peaks at 2k_F with energy ~ U ----
  % 2k_F peak marked by the dot + the "2k_F" (x) and "8.8" (y) coral ticks.  The y tick sits at the
  % MEASURED peak 8.779 and is now labelled with that value; the bare Hubbard scale U=8t is a separate
  % statement and is made in the caption, not printed on this tick.  (v3 fix: it read "~8" at 8.779.)
  \node[circle, fill=white, draw=inkPrim, line width=0.5pt, minimum size=4.8pt, inner sep=0]
        at (axis cs:1.0,8.779) {};   % = peak(q/pi=1.0) in sqw_edges.dat
\end{axis}
\end{tikzpicture}
\end{document}
